\section{Adhésion}
\label{adhesion}

\begin{enumerate}
 \item Un membre est une association qui :
  \begin{enumerate}
   \item représente la population étudiante de premier cycle en génie d'un établissement d'enseignement supérieur canadien dont le programme de génie est agréé, ou en voie d'agrément, par le Bureau canadien d'agrément des programmes de génie
   \item représente les préoccupations de la population étudiante auprès de sa faculté et de l'administration de son établissement d'enseignement supérieur
   \item a été admise comme membre de la FCEG conformément à la présente constitution.
  \end{enumerate}
 \item Plusieurs associations peuvent demander à co-représenter la population étudiante de premier cycle en génie de leur établissement lorsqu'aucune association unique ne correspond à la description.
  \begin{enumerate}
   \item La personne représentante votante doit être confirmée conjointement par la présidence de chaque association à chaque assemblée générale.
  \end{enumerate}
\end{enumerate}

\subsection{Demande d'adhésion}
\label{adhesion-demande-d-adhesion}

\begin{enumerate}
 \item Pour devenir membre, une organisation doit soumettre une demande dont le contenu est à la discrétion du CA.
 \item Le CA adoptera ensuite une résolution pour accepter ou rejeter l'admission de l'organisation, et ce au plus tard lors de l'assemblée générale suivant la soumission de la demande par l'organisation.
 \item L’assemblée générale suivante ratifiera l’adhésion de l’organisation par résolution.
\end{enumerate}
